% Figure: control volume in a contracting duct, with mass fluxes in and out.
% Shows the integral continuity bookkeeping: rho1 A1 u1 in, rho2 A2 u2 out.
\begin{figure}[htbp]
\centering
\begin{tikzpicture}[
  duct/.style={draw=engblack, very thick},
  cv/.style={draw=engred, thick, dashed},
  flux/.style={draw=engblue, -{Stealth}, thick},
  label/.style={font=\small},
]
% Contracting duct walls (upper and lower)
\draw[duct] (-4,1.6) -- (0,1.6) -- (3,0.8) -- (6,0.8);
\draw[duct] (-4,-1.6) -- (0,-1.6) -- (3,-0.8) -- (6,-0.8);

% Control volume boundary (dashed): inlet face at x=-3, outlet face at x=5
\draw[cv] (-3,-1.6) -- (-3,1.6);
\draw[cv] (5,-0.8) -- (5,0.8);
\draw[cv] (-3,1.6) -- (0,1.6) -- (3,0.8) -- (5,0.8);
\draw[cv] (-3,-1.6) -- (0,-1.6) -- (3,-0.8) -- (5,-0.8);

% Inlet flux arrows
\foreach \y in {-1.0,0,1.0} {
  \draw[flux] (-3.9,\y) -- (-3.1,\y);
}
% Outlet flux arrows (faster, shorter region, longer arrows)
\foreach \y in {-0.5,0,0.5} {
  \draw[flux] (5.1,\y) -- (6.1,\y);
}

% Inlet labels
\node[label, engblue] at (-3.6,1.35) {$\rho_1 A_1 u_1$};
\node[label, anchor=north] at (-3,-1.7) {inlet face $A_1$};
% Outlet labels
\node[label, engblue] at (6.0,1.0) {$\rho_2 A_2 u_2$};
\node[label, anchor=north] at (5,-0.95) {outlet face $A_2$};

% Outward normals
\draw[-{Latex}, engred, thick] (-3,0) -- (-3.7,0);
\node[engred, font=\scriptsize, anchor=east] at (-3.7,0.2) {$\hat{n}$};
\draw[-{Latex}, engred, thick] (5,0) -- (5.7,0);
\node[engred, font=\scriptsize, anchor=west] at (5.7,0.2) {$\hat{n}$};

% Storage label inside CV
\node[engred, font=\small] at (1.0,0) {$\displaystyle\frac{d}{dt}\!\int_{\mathcal{V}}\!\rho\,dV$};

\end{tikzpicture}
\caption[Control volume in a contracting duct]{The integral
continuity balance drawn on a fixed control volume (red dashed
boundary) spanning a contracting duct. Mass enters across the inlet
face at rate $\rho_1 A_1 u_1$ and leaves across the outlet face at
rate $\rho_2 A_2 u_2$; the outward normal $\hat{n}$ points out of the
volume on every face, so the inlet flux $\rho\vec{u}\cdot\hat{n}$ is
negative and the outlet flux positive. For steady flow the stored
mass $\int_{\mathcal{V}}\rho\,dV$ does not change, so the two face
fluxes balance: $\rho_1 A_1 u_1 = \rho_2 A_2 u_2$. The side walls
carry no flux because the velocity there is tangent to the wall. The
acceleration through the contraction is purely geometric; no force has
been invoked.}
\label{fig:vol03ch11:control-volume}
\end{figure}
